
\documentclass[11pt]{article}
\usepackage{amsmath,amssymb,amsthm,mathtools}
\usepackage[margin=1in]{geometry}
\usepackage{enumitem}
\usepackage{booktabs}

\title{Step 78: Completion Balance for Positive Weil Feature Matching}
\author{Six Birds / Formed-Layer Membrane Program}
\date{}

\newtheorem{theorem}{Theorem}[section]
\newtheorem{lemma}[theorem]{Lemma}
\newtheorem{definition}[theorem]{Definition}
\newtheorem{proposition}[theorem]{Proposition}
\newtheorem{warning}[theorem]{Warning}
\newtheorem{corollary}[theorem]{Corollary}

\newcommand{\K}{\mathsf K}
\newcommand{\A}{\mathsf A}
\newcommand{\E}{\mathsf E}
\newcommand{\pr}{\mathrm{pr}}
\newcommand{\cmp}{\mathrm{cmp}}
\newcommand{\sgn}{\mathrm{sgn}}
\newcommand{\tail}{\mathrm{tail}}
\newcommand{\pole}{\mathrm{pole}}
\newcommand{\bal}{\mathrm{bal}}
\newcommand{\diag}{\mathrm{diag}}
\newcommand{\EF}{\mathrm{EF}}

\begin{document}
\maketitle

\begin{abstract}
Step 77 showed that the signed prime autocorrelation term becomes a positive shift-difference feature only after adding a diagonal repair term.  This note isolates the corresponding completion-balance problem for the RH Weil--Douglas bridge.  The prime feature alone is not a positive representative of the signed explicit formula.  The gamma, pole/completion, and tail slots must supply a positive balance form that pays the prime diagonal and any remaining signed-form mismatch.  The central gate is a Loewner/form inequality, not a trace identity.
\end{abstract}

\section{Setup}
Let
\[
    \mathcal C^-_{\log}\subset Y^-
\]
be the anti-invariant log-side Weil core.  For \(f\in \mathcal C^-_{\log}\), let
\[
    \widetilde f(t)=\overline{f(-t)},
    \qquad h_f(a)=(f*\widetilde f)(a)=\langle \tau_a f,f\rangle .
\]
Then
\[
    \widehat h_f(\xi)=|\widehat f(\xi)|^2 .
\]

Let \(w_a\ge0\) be finite-cutoff prime-power weights, for example
\[
    w_{n,L}=\frac{\Lambda(n)}{\sqrt n}\alpha_L(\log n)^2.
\]
The signed prime autocorrelation contribution is
\[
    P_{\pr,L}(h_f)=-\sum_a w_a(h_f(a)+h_f(-a)).
\]
The positive prime shift-difference feature is
\[
    q_{\pr,L}[f]=\sum_a w_a\|\tau_a f-f\|_2^2.
\]

\begin{lemma}[Prime diagonal identity]
For every \(f\in\mathcal C^-_{\log}\),
\[
    q_{\pr,L}[f]
    = P_{\pr,L}(h_f)+d_{\pr,L}\|f\|_2^2,
\]
where
\[
    d_{\pr,L}=2\sum_a w_a.
\]
Equivalently, in Fourier variables,
\[
    \Psi_{\pr,L}(\xi)=d_{\pr,L}-2\sum_a w_a\cos(a\xi)
    =2\sum_a w_a(1-\cos(a\xi))\ge0.
\]
\end{lemma}

\begin{proof}
Expand
\[
\|\tau_a f-f\|_2^2=2\|f\|_2^2-\langle \tau_a f,f\rangle-\langle f,\tau_a f\rangle
=2h_f(0)-h_f(a)-h_f(-a).
\]
Sum with weights \(w_a\).  The Fourier expression follows from Plancherel.
\end{proof}

\section{Signed formula versus positive feature package}
The completed signed explicit-formula side may be written abstractly as
\[
    q_{\sgn,L}[f]
    =P_{\pr,L}(h_f)+R_{\sgn,L}[f],
\]
where \(R_{\sgn,L}\) collects signed archimedean, pole, completion, support, and tail terms.

The candidate positive completed carrier feature package has the form
\[
    W_{\cmp,L}=\begin{bmatrix}W_{\pr,L}\\ W_{\infty,L}\\ W_{\pole,L}\\ W_{\tail,L}\end{bmatrix},
\]
with associated positive form
\[
    q_{\cmp,L}[f]=\|W_{\cmp,L}f\|^2
    =q_{\pr,L}[f]+q_{\infty,L}^+[f]+q_{\pole,L}^+[f]+q_{\tail,L}^+[f].
\]
Thus
\[
    q_{\cmp,L}[f]-q_{\sgn,L}[f]
    =d_{\pr,L}\|f\|^2
     +q_{\infty,L}^+[f]+q_{\pole,L}^+[f]+q_{\tail,L}^+[f]
     -R_{\sgn,L}[f].
\]

\begin{definition}[Completion-balance form]
The completion-balance form is
\[
    \mathcal B_L[f]
    :=d_{\pr,L}\|f\|^2
     +q_{\infty,L}^+[f]+q_{\pole,L}^+[f]+q_{\tail,L}^+[f]
     -R_{\sgn,L}[f].
\]
The positive feature package matches the signed completed formula with no deficit when
\[
    \mathcal B_L[f]\ge0\qquad (f\in\mathcal C^-_{\log}).
\]
If only
\[
    \mathcal B_L[f]\ge -e_L[f]
\]
with \(e_L\ge0\), then \(e_L\) is the explicit-formula matching defect.
\end{definition}

\begin{theorem}[Signed-to-positive matching gate]
Assume a signed Weil domination inequality on the core,
\[
    q_Z[f]\le q_{\sgn,L}[f]+e_{\mathrm{Weil},L}[f],
\]
where \(q_Z[f]=\|V_Zf\|^2\).  If the completion-balance form satisfies
\[
    \mathcal B_L[f]\ge -e_{\bal,L}[f],
\]
then
\[
    q_Z[f]\le q_{\cmp,L}[f]+e_{\mathrm{Weil},L}[f]+e_{\bal,L}[f].
\]
In particular, if both defects vanish and the resulting form inequality extends from a form core to the completed anti-invariant response space, then
\[
    V_Z^*V_Z\preceq W_{\cmp,L}^*W_{\cmp,L},
\]
so there exists a contraction \(T\) such that
\[
    V_Z=T W_{\cmp,L},\qquad \|T\|\le1.
\]
\end{theorem}

\begin{proof}
By definition,
\[
q_{\cmp,L}=q_{\sgn,L}+\mathcal B_L.
\]
Combining the two displayed assumptions gives
\[
q_Z\le q_{\sgn,L}+e_{\mathrm{Weil},L}
\le q_{\cmp,L}+e_{\mathrm{Weil},L}+e_{\bal,L}.
\]
If the defects vanish and the core-to-full gate passes, the Loewner domination follows.  Douglas factorization gives the contraction.
\end{proof}

\section{What the gamma, pole, and tail slots must do}
The preceding theorem shows that the gamma, pole, and tail slots are not decorations.  They must pay the form
\[
    R_{\sgn,L}[f]-d_{\pr,L}\|f\|^2
\]
after adding positive completions.  Equivalently, they must make \(\mathcal B_L\) nonnegative or leave a declared defect.

\begin{proposition}[Balance split]
Suppose
\[
    R_{\sgn,L}=R_{\infty,L}^{\sgn}+R_{\pole,L}^{\sgn}+R_{\tail,L}^{\sgn}.
\]
If positive forms \(q^+_{\infty,L},q^+_{\pole,L},q^+_{\tail,L}\) satisfy
\[
    q^+_{\infty,L}-R_{\infty,L}^{\sgn}
    +q^+_{\pole,L}-R_{\pole,L}^{\sgn}
    +q^+_{\tail,L}-R_{\tail,L}^{\sgn}
    +d_{\pr,L}\|\cdot\|^2\ge0,
\]
then the signed completed carrier admits a positive feature package on the core.  If the left hand side is negative on some vector, that vector is a signed-to-positive matching witness.
\end{proposition}

\begin{warning}[Trace is not balance]
The equality
\[
    \operatorname{tr} q_{\cmp,L}=\operatorname{tr}q_{\sgn,L}
\]
for a finite presentation does not imply \(\mathcal B_L\ge0\).  Completion balance is a Loewner/form inequality, not a scalar accounting identity.
\end{warning}

\section{Candidate statuses}
The prime-power slot is already positive after diagonal repair.  The exact matching status is not earned until the diagonal repair is accounted for inside the completed formula.

\begin{center}
\begin{tabular}{lll}
\toprule
Slot & positive feature status & matching obligation\\
\midrule
Prime & shift-difference feature after diagonal repair & pay/locate \(d_{\pr,L}\|f\|^2\)\\
Gamma & plausible renormalized shift or spectral feature & prove exact archimedean matching\\
Pole/completion & finite-rank positive feature & handle pole/null modes and diagonal balance\\
Tail/support & positive feature or defect budget & prove fixed/exhaustive tail squeeze\\
\bottomrule
\end{tabular}
\end{center}

\section{Nonclaims}
This note does not prove RH.  It does not prove the classical Weil form has the required positive feature representation.  It proves only the algebraic gate: a signed explicit formula can be used for the Douglas bridge only after a positive completion-balance form is supplied, with all residuals explicitly budgeted.

\end{document}
